% f2015tjw_q10_views.tex —— 2015 天津文科卷 Q10 三视图（源 p1 栏 2）
%
% 源：.work/phase3-b10/source/2015天津文.pdf（p1 印 987）；实测清单 figure-inventory.md §F03
% 坐标换算：x_pic = x_pt − 490.0 ；y_pic = 605.0 − y_pt
%   （源页 pt、原点左上、y 向下；画布 x' 3.36–197.7、y' 6.74–195.59，单位 pt）
% 线宽：0.398 pt（源 `0.398 w` 实测）；全图实线（p1 无虚线）。
% 结构：正视图六边形 + 侧视图圆（含半径引导线）+ 俯视图六边形（= 正视图整体下移 113.388 pt）
%   ＋三组刻度＋左右竖尺寸线与横向尺寸链。
% 箭头：TikZ 标准 `>=stealth`（既定重绘约定）。源箭尖为 8 点填充多边形（实测 2.88–3.36 宽 ×
%   3.6–4.0 长），其提取坐标与墨迹不符（清单 §S-1：沿线采样仅 19.22% 落墨）；
%   本文件**不采用提取值**，24 枚中每一枚均由本次 600 dpi 墨迹独立实测（沿轴扫描取墨迹最远端）：
%     正/俯视图左右竖尺寸线 16 枚：尖分别落于对应刻度 y（实测值见下 16 行注释；与刻度的墨迹
%       间隙 0.00–0.52 pt，其中 6 枚与刻度墨迹直接相接，系尖点抗锯齿衰减，端点取刻度的几何 y）
%     横向尺寸链 6 枚（y_pt = 471.776）：尖分别为 501.860(←)、530.207(→)、530.207(←)、
%       586.900(→)、586.900(←)、615.247(→)——即**每枚尖都落在同侧刻度上**
%     侧视图半径 1 枚：尖实测 (691.00, 418.56)（距圆心 27.05，R=28.347）
%   注：600 dpi 墨迹全图（x 490–700, y 405–600）共检出**23 枚**实心箭尖（4+4 竖[正视]、4+4 竖[俯视]、
%   横向链 6、半径 1）。内容流提取计数为 24，多出的 1 枚为横向尺寸链上的「另 1 个」；
%   墨迹中不存在该枚（横向链的 6 段各只配 1 枚，两两指向同一刻度），登记为提取侧幻影。
% 存疑照源复刻：§S-7 半径箭尖距圆心 27.05 < R=28.347（箭头内缩）；§S-8 横向链 6 段长度与
%   「1|2|1」标注不成正比；§S-9 正/俯视图间距不对应「长对正、高平齐」——均按源不修正。
\begin{tikzpicture}[x=1pt, y=1pt, line width=0.398pt, >=stealth, every node/.style={font=\fontsize{10.5pt}{12.6pt}\selectfont,inner sep=0pt}]

  % ---- 正视图（六边形 + 两条内竖线） ----
  \draw (11.860,167.240) -- (40.207,138.893) -- (96.900,138.893) -- (125.247,167.240)
        -- (96.900,195.587) -- (40.207,195.587) -- cycle;
  \draw (40.207,138.893) -- (40.207,195.587);
  \draw (96.900,138.893) -- (96.900,195.587);

  % ---- 侧视图（圆 + 半径引导线，箭尖＝墨迹实测 (691.00,418.56)） ----
  \draw (181.941,167.240) circle [radius=28.347];
  \draw[->] (181.941,167.240) -- (201.000,186.440);

  % ---- 俯视图（同上六边形整体下移 113.388 pt） ----
  \draw (11.860,53.852) -- (40.207,25.505) -- (96.900,25.505) -- (125.247,53.852)
        -- (96.900,82.199) -- (40.207,82.199) -- cycle;
  \draw (40.207,25.505) -- (40.207,82.199);
  \draw (96.900,25.505) -- (96.900,82.199);

  % ---- 刻度：正视图左/右各 3、顶边 4；俯视图左/右各 3 ----
  \draw (3.356,138.893) -- (9.025,138.893);    \draw (128.082,138.893) -- (133.751,138.893);
  \draw (3.356,167.240) -- (9.025,167.240);    \draw (128.082,167.240) -- (133.751,167.240);
  \draw (3.356,195.587) -- (9.025,195.587);    \draw (128.082,195.587) -- (133.751,195.587);
  \draw (11.860,130.389) -- (11.860,136.058);  \draw (40.207,130.389) -- (40.207,136.058);
  \draw (96.900,130.389) -- (96.900,136.058);  \draw (125.247,130.389) -- (125.247,136.058);
  \draw (3.356,25.506) -- (9.025,25.506);      \draw (128.082,25.506) -- (133.751,25.506);
  \draw (3.356,53.852) -- (9.025,53.852);      \draw (128.082,53.852) -- (133.751,53.852);
  \draw (3.356,82.199) -- (9.025,82.199);      \draw (128.082,82.199) -- (133.751,82.199);

  % ---- 正视图左右竖尺寸线（各 4 段，箭尖落于刻度；括号内为 600dpi 墨迹实测尖 y_pt） ----
  \draw[->] (130.917,147.397) -- (130.917,138.893);   % 457.603 -> 尖 465.96, 刻度 466.107
  \draw[->] (130.917,158.736) -- (130.917,167.240);   % 446.264 -> 尖 438.00, 刻度 437.760
  \draw[->] (130.917,175.744) -- (130.917,167.240);   % 429.256 -> 尖 437.28, 刻度 437.760
  \draw[->] (130.917,187.083) -- (130.917,195.587);   % 417.917 -> 尖 409.44, 刻度 409.413
  \draw[->] (  6.191,147.397) -- (  6.191,138.893);   % 同上（左，实测 465.96）
  \draw[->] (  6.191,158.736) -- (  6.191,167.240);   % 同上（左，实测 438.00）
  \draw[->] (  6.191,175.744) -- (  6.191,167.240);   % 同上（左，实测 437.28）
  \draw[->] (  6.191,187.083) -- (  6.191,195.587);   % 同上（左，实测 409.44）

  % ---- 俯视图左右竖尺寸线（各 4 段） ----
  \draw[->] (130.917,34.010) -- (130.917,25.506);     % 570.990 -> 尖 579.00, 刻度 579.494
  \draw[->] (130.917,45.348) -- (130.917,53.852);     % 559.652 -> 尖 551.40, 刻度 551.148
  \draw[->] (130.917,62.357) -- (130.917,53.852);     % 542.643 -> 尖 550.92, 刻度 551.148
  \draw[->] (130.917,73.695) -- (130.917,82.199);     % 531.305 -> 尖 523.32, 刻度 522.801
  \draw[->] (  6.191,34.010) -- (  6.191,25.506);     % 同上（左，实测 579.00）
  \draw[->] (  6.191,45.348) -- (  6.191,53.852);     % 同上（左，实测 551.40）
  \draw[->] (  6.191,62.357) -- (  6.191,53.852);     % 同上（左，实测 550.92）
  \draw[->] (  6.191,73.695) -- (  6.191,82.199);     % 同上（左，实测 523.32）

  % ---- 横向尺寸链（6 段，箭尖均落在刻度上） ----
  \draw[->] (20.364,133.224) -- (11.860,133.224);     % 尖 501.860(←)
  \draw[->] (31.703,133.224) -- (40.207,133.224);     % 尖 530.207(→)
  \draw[->] (62.884,133.224) -- (40.207,133.224);     % 尖 530.207(←)
  \draw[->] (74.223,133.224) -- (96.900,133.224);     % 尖 586.900(→)
  \draw[->] (105.404,133.224) -- (96.900,133.224);    % 尖 586.900(←)
  \draw[->] (116.743,133.224) -- (125.247,133.224);   % 尖 615.247(→)

  % ---- 标签（视图名：源 pdftotext -bbox 字形盒左上角；尺寸数字：源字形盒左上角，
  %      竖线数字与半径数字的 x 由 600dpi 墨迹反推（清单只给 y 档）——inner sep=0） ----
  \node[anchor=north west,inner sep=0] at ( 52.80,122.06) {正视图};
  \node[anchor=north west,inner sep=0] at (166.19,122.07) {侧视图};
  \node[anchor=north west,inner sep=0] at ( 52.80, 17.24) {俯视图};
  \node[anchor=north west,inner sep=0] at ( 23.41,137.13) {$1$};
  \node[anchor=north west,inner sep=0] at ( 65.93,137.13) {$2$};
  \node[anchor=north west,inner sep=0] at (108.45,137.13) {$1$};
  \node[anchor=north west,inner sep=0] at (193.41,173.98) {$1$};
  \node[anchor=north west,inner sep=0] at (  3.55,185.32) {$1$};
  \node[anchor=north west,inner sep=0] at (  3.55,156.97) {$1$};
  \node[anchor=north west,inner sep=0] at (  3.55, 71.93) {$1$};
  \node[anchor=north west,inner sep=0] at (  3.55, 43.58) {$1$};
  \node[anchor=north west,inner sep=0] at (128.31,185.32) {$1$};
  \node[anchor=north west,inner sep=0] at (128.31,156.97) {$1$};
  \node[anchor=north west,inner sep=0] at (128.31, 71.93) {$1$};
  \node[anchor=north west,inner sep=0] at (128.31, 43.58) {$1$};

\end{tikzpicture}
